\documentclass[11pt,a4paper]{scrreprt}

% ---------- Layout ----------
\usepackage[left=2.2cm,right=2.2cm,top=2.6cm,bottom=2.8cm,headheight=14pt]{geometry}
\usepackage[utf8]{inputenc}
\usepackage[T1]{fontenc}
\usepackage{lmodern}
\usepackage{microtype}
\usepackage{ragged2e}

% ---------- Color & tables ----------
\usepackage[table]{xcolor}
\usepackage{colortbl}
\usepackage{booktabs}
\usepackage{array}
\usepackage{tabularx}
\usepackage{longtable}

% ---------- Math ----------
\usepackage{amsmath}
\usepackage{amssymb}

% ---------- Lists, boxes, graphics ----------
\usepackage{enumitem}
\usepackage{tcolorbox}
\tcbuselibrary{skins}
\usepackage{tikz}
\usetikzlibrary{positioning,calc,shapes.geometric,arrows.meta,backgrounds}
\usepackage{fontawesome5}
\usepackage{titlesec}
\usepackage{scrlayer-scrpage}
\usepackage[hidelinks]{hyperref}

% ---------- Palette ----------
\definecolor{sopnavy}{RGB}{17,42,66}
\definecolor{sopteal}{RGB}{0,116,108}
\definecolor{sopgray}{RGB}{92,99,108}
\definecolor{sopamber}{RGB}{178,122,17}
\definecolor{soplight}{RGB}{241,244,247}
\definecolor{sopline}{RGB}{206,214,221}

% ---------- Heading style ----------
\titleformat{\section}
  {\color{sopnavy}\Large\bfseries}
  {\thesection}{0.5em}{}
  [{\color{sopline}\titlerule[0.9pt]}]
\titlespacing*{\section}{0pt}{1.4em}{0.8em}

\titleformat{\subsection}
  {\color{sopteal}\large\bfseries}
  {\thesubsection}{0.5em}{}
\titlespacing*{\subsection}{0pt}{1.0em}{0.5em}

\titleformat{\subsubsection}
  {\color{sopgray}\normalsize\bfseries\itshape}
  {\thesubsubsection}{0.5em}{}
\titlespacing*{\subsubsection}{0pt}{0.8em}{0.3em}

% ---------- Header / Footer ----------
\pagestyle{scrheadings}
\clearpairofpagestyles
\ihead{\small\color{sopgray}SOP-CT-014 \textbar\ Multicenter Clinical Trial Conduct}
\ohead{\small\color{sopgray}Version 4.0}
\ifoot{\small\color{sopgray}Meridian Clinical Research Network -- Confidential}
\ofoot{\small\color{sopgray}Page \thepage\ of \pageref{LastPage}}
\setkomafont{pageheadfoot}{\normalfont}
\usepackage{lastpage}

\setlist[itemize]{leftmargin=1.4em,itemsep=2pt,topsep=2pt}
\setlist[enumerate]{leftmargin=1.6em,itemsep=2pt,topsep=2pt}

\newcommand{\gcpref}[1]{{\color{sopteal}\small\itshape #1}}

\begin{document}

% ======================================================================
% TITLE PAGE
% ======================================================================
\begin{titlepage}
\thispagestyle{empty}
\begin{tikzpicture}[remember picture,overlay]
  \fill[sopnavy] (current page.north west) rectangle
      ([yshift=-4.6cm]current page.north east);
  \fill[sopteal] (current page.north west) rectangle
      ([xshift=0.28cm,yshift=-4.6cm]current page.north west);
  \node[anchor=north west,text=white,font=\small\bfseries] at
      ([xshift=2.2cm,yshift=-1.1cm]current page.north west)
      {MERIDIAN CLINICAL RESEARCH NETWORK};
  \node[anchor=north east,text=white!85!black,font=\small] at
      ([xshift=-2.2cm,yshift=-1.1cm]current page.north east)
      {\faLock\ CONTROLLED DOCUMENT};
  \node[anchor=west,text=white,font=\Huge\bfseries] at
      ([xshift=2.2cm,yshift=-2.7cm]current page.north west)
      {Standard Operating Procedure};
  \node[anchor=west,text=white!80!sopteal,font=\Large] at
      ([xshift=2.2cm,yshift=-3.7cm]current page.north west)
      {Conduct of Multicenter Clinical Trials -- GCP Compliance};
\end{tikzpicture}

\vspace*{5.1cm}

\begin{center}
{\color{sopgray}\large Protocol NX-2214-301}\\[2pt]
{\small\itshape A Phase III, Randomized, Double-Blind, Placebo-Controlled Multicenter\\
Study Evaluating the Efficacy and Safety of NX-2214 in Adults with\\
Moderate-to-Severe Plaque Psoriasis}
\end{center}

\vspace{0.9cm}

\noindent
\begin{tabularx}{\linewidth}{@{}l X@{}}
\rowcolor{soplight}
\textbf{SOP Number} & SOP-CT-014 \\
\textbf{Version} & 4.0 \\
\rowcolor{soplight}
\textbf{Effective Date} & 01 June 2026 \\
\textbf{Supersedes} & Version 3.2 (12 November 2024) \\
\rowcolor{soplight}
\textbf{Next Scheduled Review} & 01 June 2027 \\
\textbf{Sponsor} & Nexa Therapeutics, Inc. \\
\rowcolor{soplight}
\textbf{Coordinating Center} & Vertex Clinical Trials Unit \\
\textbf{Governing CRO} & Meridian Clinical Research Network \\
\end{tabularx}

\vspace{1.0cm}
{\color{sopnavy}\bfseries Approval Signatures}\\[4pt]
\noindent
\begin{tabularx}{\linewidth}{@{}>{\raggedright\arraybackslash}p{3.6cm} >{\raggedright\arraybackslash}p{4.4cm} X >{\centering\arraybackslash}p{2.1cm}@{}}
\toprule
\textbf{Role} & \textbf{Name} & \textbf{Signature} & \textbf{Date} \\
\midrule
Sponsor Medical Director & Dr. Elena Vasquez, Nexa Therapeutics & \rule{3.2cm}{0.4pt} & \rule{1.6cm}{0.4pt} \\
Coordinating Investigator & Dr. Michael Chen, Vertex Clinical Trials Unit & \rule{3.2cm}{0.4pt} & \rule{1.6cm}{0.4pt} \\
Head of Clinical Operations & Sarah Okonkwo, Meridian CRN & \rule{3.2cm}{0.4pt} & \rule{1.6cm}{0.4pt} \\
QA/QC Director & Dr. Priya Nair, Meridian CRN & \rule{3.2cm}{0.4pt} & \rule{1.6cm}{0.4pt} \\
\bottomrule
\end{tabularx}

\vfill
{\footnotesize\color{sopgray}\raggedright
This document is proprietary to Nexa Therapeutics, Inc. and Meridian Clinical Research
Network. Unauthorized reproduction or distribution outside the study team is prohibited.
Uncontrolled when printed -- verify current version in the electronic Trial Master File (eTMF)
prior to use.\par}
\end{titlepage}

% ======================================================================
% CHAPTER-LEVEL CONTENT (flat, no forced page breaks)
% ======================================================================

\section{Purpose and Scope}

This Standard Operating Procedure (SOP) defines the mandatory processes for planning,
initiating, conducting, monitoring, and closing out multicenter interventional clinical trials
sponsored by Nexa Therapeutics, Inc. and operationally managed by Meridian Clinical Research
Network. It establishes uniform expectations for all personnel -- sponsor staff, contract
research organization (CRO) staff, Principal Investigators, and site personnel -- engaged in
Protocol NX-2214-301 and any subsequent multicenter protocol referencing this SOP.

This procedure applies to all 42 participating investigator sites across 9 countries and covers
the full trial lifecycle from site feasibility assessment through database lock and archiving.
It does not supersede local regulatory requirements, institutional policies, or the approved
protocol; where a conflict exists, the more stringent requirement governs.

\begin{tcolorbox}[
  colback=soplight, colframe=sopnavy, boxrule=0.6pt, arc=2pt,
  left=8pt, right=8pt, top=6pt, bottom=6pt,
  title={\faFileMedical\ \ Applicable Regulations and Guidance},
  coltitle=white, colbacktitle=sopnavy, fonttitle=\bfseries\small,
  attach boxed title to top left={xshift=8pt,yshift=-2pt},
  boxed title style={boxrule=0pt,arc=2pt}
]
\small
ICH E6(R2) Good Clinical Practice \textbullet\ 21 CFR Parts 50, 54, 56 and 312 (US FDA)
\textbullet\ Regulation (EU) No 536/2014 on Clinical Trials \textbullet\ Declaration of
Helsinki (2013, Fortaleza revision) \textbullet\ Applicable local data protection and health
privacy law at each participating site.
\end{tcolorbox}

\section{Roles and Responsibilities}

Each role below carries defined accountability under this SOP. Delegation of tasks at site
level must be recorded on the study-specific Delegation of Authority (DoA) log and is limited
to individuals with documented training and qualifications.

\begin{center}
\small
\begin{tabularx}{\linewidth}{@{} >{\raggedright\arraybackslash}p{3.3cm} X >{\raggedright\arraybackslash}p{3.0cm} @{}}
\toprule
\rowcolor{sopnavy}
{\color{white}\textbf{Role}} & {\color{white}\textbf{Key Responsibilities}} & {\color{white}\textbf{Reports To}} \\
\midrule
Sponsor (Nexa Therapeutics) & Overall trial accountability, protocol authorship, safety
  determinations, regulatory submissions & Board of Directors \\
\rowcolor{soplight}
Coordinating Investigator & Scientific leadership across sites, protocol interpretation
  queries, publication oversight & Sponsor \\
Principal Investigator (per site) & Subject safety, informed consent, protocol adherence,
  supervision of delegated staff & Coordinating Investigator \\
\rowcolor{soplight}
Clinical Research Associate & Site monitoring, source data verification, regulatory document
  maintenance & Clinical Operations Lead \\
Data Management Team & eCRF design, query resolution, database lock readiness & Head of
  Clinical Operations \\
\rowcolor{soplight}
Pharmacovigilance Unit & SAE triage, expedited reporting, aggregate safety review & Sponsor
  Medical Director \\
IRB / Independent Ethics Committee & Protocol and consent approval, ongoing ethical oversight
  & Institution (independent) \\
\rowcolor{soplight}
Site Study Coordinator & Visit scheduling, consent logistics, IP dispensing records &
  Principal Investigator \\
\bottomrule
\end{tabularx}
\end{center}

\section{Site Selection, Initiation, and Activation}

Sites are selected based on feasibility survey results, prior trial performance, patient
population access, and infrastructure adequacy (pharmacy, biosample handling, emergency
equipment). Activation follows a fixed sequence; no site may enroll subjects prior to
completion of every step below.

\begin{center}
\begin{tikzpicture}[
  node distance=4mm and 3mm,
  box/.style={rectangle, rounded corners=3pt, draw=sopnavy, fill=soplight,
    text width=2.15cm, align=center, minimum height=1.15cm, font=\scriptsize, inner sep=3pt},
  arrow/.style={-{Stealth[length=2mm]}, sopteal, thick}
]
\node[box] (n1) {Feasibility Assessment};
\node[box, right=of n1] (n2) {Site Qualification Visit};
\node[box, right=of n2] (n3) {Contract \& Budget Execution};
\node[box, right=of n3] (n4) {IRB/EC Approval};
\node[box, right=of n4] (n5) {Site Initiation Visit};
\node[box, right=of n5, fill=sopteal, text=white, draw=sopnavy] (n6) {Activation -- First Subject In};
\draw[arrow] (n1) -- (n2);
\draw[arrow] (n2) -- (n3);
\draw[arrow] (n3) -- (n4);
\draw[arrow] (n4) -- (n5);
\draw[arrow] (n5) -- (n6);
\end{tikzpicture}
\end{center}

The target interval from feasibility confirmation to activation is 45 calendar days. Sites
exceeding 75 days without activation are escalated to the Steering Committee for a
go/no-go re-evaluation. The Site Initiation Visit (SIV) must cover protocol training, IP
handling, electronic data capture (EDC) system access, and safety reporting pathways, and
must be documented on the SIV attendance log retained in the Trial Master File (TMF).

\section{Informed Consent Process}

Informed consent is obtained by the Investigator or a qualified, delegated sub-investigator
prior to any protocol-specific procedure. The consent discussion must occur in a private
setting, in the subject's preferred language, and allow adequate time for questions.

\begin{minipage}[t]{0.485\linewidth}
\begin{tcolorbox}[colback=white, colframe=sopteal, boxrule=0.7pt, arc=2pt,
  title={\faCheckCircle\ Mandatory Elements}, coltitle=white, colbacktitle=sopteal,
  fonttitle=\bfseries\small, left=6pt, right=6pt, top=4pt, bottom=4pt]
\small
\begin{itemize}
\item Current IRB/EC-approved consent version
\item Voluntary participation and right to withdraw
\item Foreseeable risks and expected benefits
\item Alternative treatments available
\item Confidentiality and data-sharing scope
\end{itemize}
\end{tcolorbox}
\end{minipage}\hfill
\begin{minipage}[t]{0.485\linewidth}
\begin{tcolorbox}[colback=white, colframe=sopamber, boxrule=0.7pt, arc=2pt,
  title={\faExclamationTriangle\ Re-Consent Triggers}, coltitle=white, colbacktitle=sopamber,
  fonttitle=\bfseries\small, left=6pt, right=6pt, top=4pt, bottom=4pt]
\small
\begin{itemize}
\item Substantive protocol amendment
\item New safety information affecting risk profile
\item Change to biological sample use or retention
\item Extension of study duration
\end{itemize}
\end{tcolorbox}
\end{minipage}
\par

Signed and dated consent forms are filed in the subject's source record; the process must be
documented contemporaneously in the source notes, including the date, time, and version of
the consent form used. Non-English-speaking subjects must receive a certified translated
consent form -- verbal translation alone is not acceptable except under emergency exception
provisions defined in the protocol.

\section{Investigational Product Management}

\subsection{Storage and Temperature Monitoring}

Investigational product (IP) NX-2214 and matching placebo must be stored at
2\textdegree C--8\textdegree C in a dedicated, access-controlled pharmacy refrigerator with
continuous calibrated temperature logging. Logs are reviewed by site pharmacy staff daily and
by the CRA at every monitoring visit.

\subsection{Accountability and Reconciliation}

Drug accountability logs must record every receipt, dispensing, return, and destruction event,
reconciled against subject dosing records at each visit. Discrepancies exceeding one dosing
unit must be investigated and documented within 5 business days.

\begin{tcolorbox}[colback=sopamber!8!white, colframe=sopamber, boxrule=0.8pt, arc=2pt,
  left=8pt, right=8pt, top=6pt, bottom=6pt,
  title={\faExclamationTriangle\ \ Temperature Excursion Handling}, coltitle=white,
  colbacktitle=sopamber, fonttitle=\bfseries\small]
\small
Upon any excursion outside 2\textdegree C--8\textdegree C: (1) quarantine affected product
immediately and label \emph{DO NOT USE}; (2) record duration and temperature range from the
logger; (3) notify the CRA and Sponsor Pharmacovigilance within 24 hours; (4) do not dispense
or destroy quarantined product until written disposition is received from the Sponsor.
\end{tcolorbox}

\section{Safety Reporting}

All adverse events are recorded in the eCRF regardless of causality. Serious adverse events
(SAEs) follow expedited timelines independent of routine visit schedules.

\begin{center}
\small
\begin{tabularx}{\linewidth}{@{} >{\raggedright\arraybackslash}p{4.4cm} >{\raggedright\arraybackslash}p{3.3cm} X @{}}
\toprule
\rowcolor{sopnavy}
{\color{white}\textbf{Event Type}} & {\color{white}\textbf{Reporting Timeline}} &
{\color{white}\textbf{Recipient}} \\
\midrule
SAE -- fatal or life-threatening & Within 24 hours of awareness & Sponsor Pharmacovigilance,
  IRB/EC \\
\rowcolor{soplight}
SAE -- other serious & Within 24 hours of awareness & Sponsor Pharmacovigilance \\
Suspected Unexpected Serious Adverse Reaction (SUSAR) & Within 7 or 15 calendar days per
  expectedness & Regulatory Authorities, all sites \\
\rowcolor{soplight}
Pregnancy during study exposure & Within 24 hours of awareness & Sponsor Pharmacovigilance \\
Overdose (with or without associated AE) & Within 24 hours of awareness & Sponsor
  Pharmacovigilance \\
\bottomrule
\end{tabularx}
\end{center}

\gcpref{Reference: ICH E2A, Clinical Safety Data Management -- Definitions and Standards for
Expedited Reporting.}

\section{Monitoring and Quality Oversight}

Monitoring combines on-site visits with centralized statistical data review to detect outliers,
site drift, and data integrity concerns between visits.

\begin{center}
\small
\begin{tabularx}{\linewidth}{@{} >{\raggedright\arraybackslash}p{3.9cm} >{\raggedright\arraybackslash}p{2.8cm} X @{}}
\toprule
\rowcolor{sopteal}
{\color{white}\textbf{Visit Type}} & {\color{white}\textbf{Frequency}} &
{\color{white}\textbf{Primary Focus}} \\
\midrule
Site Initiation Visit (SIV) & Once, pre-activation & Training, readiness, regulatory binder
  completeness \\
\rowcolor{soplight}
Interim Monitoring Visit (IMV) & Every 6--8 weeks & Source data verification, consent
  compliance, IP accountability \\
Remote / Centralized Monitoring & Continuous & Query trends, enrollment pace, safety signal
  detection \\
\rowcolor{soplight}
Close-Out Visit (COV) & Once, post-database lock & Reconciliation, IP return, archiving
  confirmation \\
For-Cause Audit & As triggered & Root-cause investigation of significant findings \\
\bottomrule
\end{tabularx}
\end{center}

\section{Protocol Deviations and CAPA}

Deviations are classified as \textbf{Minor} (no impact on subject safety or data integrity),
\textbf{Major} (potential impact requiring corrective action), or \textbf{Critical}
(direct impact on subject safety, rights, or primary endpoint reliability).

\begin{tcolorbox}[colback=soplight, colframe=sopnavy, boxrule=0.6pt, arc=2pt,
  left=8pt, right=8pt, top=6pt, bottom=6pt,
  title={\faClipboardList\ \ Reporting and CAPA Requirement}, coltitle=white,
  colbacktitle=sopnavy, fonttitle=\bfseries\small]
\small
Major and Critical deviations must be reported to the CRA within 3 business days and to the
IRB/EC per local requirements. A Corrective and Preventive Action (CAPA) plan is mandatory
for all Major and Critical deviations, with root-cause analysis, assigned owner, and a
verification-of-effectiveness date logged in the Quality Issue Tracker.
\end{tcolorbox}

\section{Records Retention and Archiving}

Essential documents, including signed consent forms, source records, monitoring visit reports,
and the completed regulatory binder, are retained for a minimum of 15 years following trial
completion, or longer where local law requires. Site archiving is confirmed at the Close-Out
Visit and documented in the TMF Index prior to final site release.

\vspace{0.6cm}
\noindent{\color{sopline}\rule{\linewidth}{0.8pt}}
\vspace{0.4cm}

\section*{Appendix A -- Revision History}
\addcontentsline{toc}{section}{Appendix A -- Revision History}

\begin{center}
\small
\begin{tabularx}{\linewidth}{@{} >{\centering\arraybackslash}p{1.3cm} >{\centering\arraybackslash}p{2.3cm} X @{}}
\toprule
\rowcolor{sopgray}
{\color{white}\textbf{Version}} & {\color{white}\textbf{Date}} & {\color{white}\textbf{Summary of Changes}} \\
\midrule
1.0 & 04 Feb 2022 & Initial issue for protocol NX-2214-101 (Phase I). \\
\rowcolor{soplight}
2.0 & 19 Sep 2022 & Expanded to multicenter scope; added CRA monitoring cadence table. \\
3.0 & 22 Jan 2024 & Added centralized/remote monitoring; revised SAE timelines to align
  with amended ICH E2A guidance. \\
\rowcolor{soplight}
3.2 & 12 Nov 2024 & Clarified temperature excursion quarantine procedure; added pregnancy
  reporting pathway. \\
4.0 & 01 Jun 2026 & Full rewrite for Protocol NX-2214-301; added deviation classification
  matrix and CAPA requirement; updated retention period to 15 years. \\
\bottomrule
\end{tabularx}
\end{center}

\section*{Appendix B -- Regulatory References}
\addcontentsline{toc}{section}{Appendix B -- Regulatory References}

\small
\begin{itemize}
\item ICH E6(R2) -- Good Clinical Practice: Integrated Addendum
\item ICH E2A -- Clinical Safety Data Management: Definitions and Standards for Expedited
  Reporting
\item US Code of Federal Regulations, Title 21, Parts 50, 54, 56, and 312
\item Regulation (EU) No 536/2014 on Clinical Trials on Medicinal Products for Human Use
\item World Medical Association Declaration of Helsinki (2013, Fortaleza revision)
\end{itemize}

\vspace{0.8cm}
\noindent{\footnotesize\color{sopgray}\faLock\ This SOP is uncontrolled when printed. Refer to
the Meridian eQMS for the current approved version before use. \textbullet\ Document owner:
Head of Clinical Operations, Meridian Clinical Research Network.}

\end{document}
